\documentclass[11pt]{article}
\usepackage[margin=1in]{geometry}
\usepackage{amsmath,amssymb,amsthm}
\newtheorem{theorem}{Theorem}
\newtheorem{lemma}{Lemma}
\newtheorem{corollary}{Corollary}
\newcommand{\E}{\mathbb E}
\newcommand{\Var}{\operatorname{Var}}
\newcommand{\su}{\mathfrak{su}}
\title{A two-link obstruction to removing the physical H contact}\author{}\date{}
\begin{document}\maketitle
\noindent\textbf{Scope.} This is a finite-cutoff, nonlinear, same-Wilson result.
It strengthens the manuscript's ultraviolet lower bound and proves a lower
bound for its specified ordinary remainder. It neither proves a matching
upper bound nor closes all-history MT36. No source prescription is altered.


\section{Precise statement and conventions}
Use AI1--3 (PDF pp.~635--636), MT1 (p.~681), and PD12.20--21 (p.~1186).
Let $G=SU(N)$, $N\ge2$, with $-\operatorname{tr}(T^AT^B)=\delta_{AB}$.
On the periodic four-dimensional lattice of spacing $a$, side $T=na$, put
\[
 s_p=N-\Re\operatorname{tr}U_p,\qquad
 S_0=\sum_p s_p,\qquad
 d\mu_\beta=Z_\beta^{-1}e^{-\beta S_0}\prod_b dU_b,
 \quad \beta=c_\beta g^{-2},\quad c_\beta>0.
\]
Haar factors are normalized. Bond derivatives are exactly
$X_b^AF=\left.\frac d{dt}F(\ldots,e^{tT^A}U_b,\ldots)\right|_{0}$.
Write $i=(b,A)$, with the positive bond direction denoted $\nu_b$, and
\[
 F_i=X_iS_0,\quad F_{i;\mu}=X_iS_0^{(\mu\nu_b)},\quad
 D[f]=\sum_i f_i\sum_{\mu\ne\nu_b}F_{i;\mu}^2,\quad
 P[f]=\sum_i f_iF_i^2,\quad \ell[f]=\sum_i f_iX_i^2S_0.
\]
All incidences, color sums and transformed tests in the fixed pointed
completion are retained. Set
\[
 G_H[f]=D[f]-P[f]/3,\qquad \sigma=g^{-2}a^{-2},\qquad
 O_H(f)=\sigma G_H[f],\qquad
 r_H(f)=\sigma\{D[f]-\ell[f]/(3\beta)\}.
\]
For $f=1$ these global sums are already invariant under the fixed pointed
hypercubic averaging. The exact Casimir formula is
\begin{equation}\label{casimir}
 \ell[f]=C_F\sum_p f_p(N-s_p),\qquad f_p=\sum_{b\in\partial p}f_b,
\end{equation}
with the corresponding finite average for a completed test. In particular,
$\ell[1]=4C_F(N\#p-S_0)$, with its constant intact.

A \emph{single-plaquette subtraction} is any real random variable
\begin{equation}\label{lower}
 L(U)=c+\sum_p h_p(U_p),
\end{equation}
where $c$ and the functions $h_p$ are deterministic, may depend arbitrarily
on $a,\beta,T,p$, and every summand belongs to $L^2(\mu_\beta)$. Smooth real class functions
are sufficient; the proof also permits measurable functions when square
integrable. No uniform coefficient or derivative bounds are imposed.
This includes arbitrary weighted Wilson action terms, constants, and the
literal Casimir companion. It does not include arbitrary products of
functions of two distinct plaquettes or the full unspecified lower sector.

\begin{theorem}[Ultraviolet fluctuation survives single-plaquette subtractions]
\label{main}
There are $c>0$, $\beta_0$ and $n_0$, depending only on the group,
$c_\beta$ and fixed lattice/pointing geometry, such that for $n\ge n_0$,
$\beta\ge\beta_0$, and every $L$ of the form \eqref{lower},
\begin{align}
 \Var_{\mu_\beta}\bigl(O_H(1)-L\bigr)&\ge cT^4a^{-8},\label{mainH}\\
 \Var_{\mu_\beta}\bigl(\sigma D[1]-L\bigr)&\ge cT^4a^{-8}.\label{mainD}
\end{align}
The constant is independent of the functions and sizes of the subtraction.
For the original fixed real physical test $f$ equal to one on a
positive-volume box, both estimates hold with $c_fa^{-8}$ for all sufficiently
small $a$, using the exact completed tests.
\end{theorem}

\begin{corollary}[The manuscript's ordinary H remainder also diverges]
Under the same hypotheses,
\[
 \Var_{\mu_\beta}(r_H(1))\ge cT^4a^{-8},\qquad
 \Var_{\mu_\beta}(r_H(f))\ge c_fa^{-8}.
\]
The assertions remain valid after subtracting any additional
single-plaquette term from $r_H$.
\end{corollary}
\begin{proof}
Take $L=\sigma\ell[f]/(3\beta)$ in \eqref{mainD}, using \eqref{casimir}.
This keeps the exact constant, original tests, and all plaquette incidences.
The sum of this term and any further subtraction is still \eqref{lower}.
\end{proof}

\section{The elementary two-link calculation}
Choose distinct coordinate directions $\nu,\mu$, and free bonds
\[
 b_-=(x-e_\mu,\nu),\qquad b_+=(x+e_\mu,\nu),
 \qquad b_0=(x,\nu).
\]
Here the offsets are in lattice units. No plaquette contains both free bonds.
Their plaquette stars have exactly one bond in common, $b_0$; the two
plaquettes meeting there are in the same $\mu\nu$ plane. All statements are
inside an embedded fixed box, so there is no periodic wrapping.

Put every other link of their determining neighborhood at the identity and
write $U_-=e^u$, $U_+=e^v$, with $u,v\in\su(N)$. In the AI1 norm,
$s(e^u)=|u|^2/2+O(|u|^3)$. The quadratic part of each plaquette action
is one half the squared oriented sum of its link coordinates. Consequently
the linearized force is the plaquette incidence matrix squared, plane by
plane. For one free link, its own three plane forces are $2u$ and eighteen
other bond forces have magnitude $|u|$. Thus $D$ has diagonal coefficient
$3\cdot4+18=30$, while $P$ has $6^2+18=54$.

For the cross term, only the force row at $b_0$ depends on both variables.
Its $\mu$-plane force is $-u-v$ to first order. Its other plane forces vanish
to first order in these variables. Thus its cross coefficient is
$2\langle u,v\rangle$ in both $D$ and $P$. No other force row contributes
a cross term. Therefore the complete quadratic jets are
\begin{align}
 D^{(2)}(u,v)&=30(|u|^2+|v|^2)+2\langle u,v\rangle,\label{jetD}\\
 P^{(2)}(u,v)&=54(|u|^2+|v|^2)+2\langle u,v\rangle,\\
 G_H^{(2)}(u,v)&=12(|u|^2+|v|^2)+\tfrac43\langle u,v\rangle.\label{jetH}
\end{align}
In particular for any fixed unit generator coordinate $A$,
\[
 \partial_{u^A}\partial_{v^A}D(0,0)=2,\qquad
 \partial_{u^A}\partial_{v^A}G_H(0,0)=4/3.
\]
These are jets of exact nonlinear observables, not a replacement of the
measure by a Gaussian. The finite calculation
independently enumerates oriented plaquettes and checks the three quadratic
matrices by integer/rational arithmetic. The argument above proves why the
enumerated coefficients are the nonlinear jets for every $SU(N)$.

By smoothness of finitely many compact plaquette words, these mixed
derivatives are bounded below by a fixed $\kappa>0$ when all bath links
in that neighborhood are sufficiently near the identity, and both free
coordinates are in a sufficiently small fixed identity chart. This uses
only a finite-dimensional neighborhood, with constants independent of
$n,a,\beta$. Gauge transformations preserving the bath-smallness description
preserve $D,G_H$ and Haar measure.

\section{A conditional variance lemma immune to additive terms}
\begin{lemma}\label{mixed}
Let $(z,w)$ have a density $p$ bounded below by $m_0>0$ on a fixed product
of balls $B\times B\subset\mathbb R^d\times\mathbb R^d$. Suppose a smooth
function $Q$ satisfies $\partial_{z^A}\partial_{w^A}Q\ge\kappa>0$ there.
For arbitrary square-integrable additive terms $A(z)+B(w)$,
\[
 \Var\bigl(Q(z,w)-A(z)-B(w)\bigr)\ge
 \frac{m_0\kappa^2}{\|\partial_{z^A}\partial_{w^A}\phi\|_{L^2}^2}>0,
\]
where any fixed $\phi\in C_c^\infty(B\times B)$, $\phi\ge0$,
$\int\phi=1$, may be used. No independence is needed for this lemma.
\end{lemma}
\begin{proof}
Let $\psi=\partial_{z^A}\partial_{w^A}\phi$. Integration by parts twice
and Fubini give $\int Q\psi\ge\kappa$, whereas $\int(A+B+c)\psi=0$ for
every constant $c$. Hence Cauchy--Schwarz yields
$\|Q-A-B-c\|_{L^2(B\times B)}\ge\kappa/\|\psi\|_2$.
Multiply its square by the density lower bound and minimize over $c$.
Locally integrable additive functions suffice, so no derivatives of the
subtractions have been assumed or estimated.
\end{proof}

\section{Uniform nonlinear conditional laws and a positive density of blocks}
Here are the same-measure details needed to use the lemma; compare the
one-link argument PD12.41--53, PDF pp.~1190--1193.

First the normalized Haar partition function obeys
\[
 -\log Z_\beta\le C\#E\log(2+\beta).
\]
Indeed restrict every link to a ball of radius $c\beta^{-1/2}$: its Haar
volume is bounded below by $c'\beta^{-\dim G/2}$, and every plaquette
energy is $O(\beta^{-1})$. Since $\E_\beta S_0$ is decreasing,
\[
 \frac\beta2\E_\beta S_0
 \le\log Z_{\beta/2}-\log Z_\beta\le-\log Z_\beta.
\]
Translation and coordinate symmetries give
$\E_\beta s_p\le C\log(2+\beta)/\beta$ uniformly in $n$.
For any fixed contractible cube $Q$, Markov and a union bound imply
\[
 \mu_\beta\{s_p\le\epsilon^2\text{ for every }p\subset Q\}
 \ge1-C_Q\epsilon^{-2}\log(2+\beta)/\beta.
\]
On that event, a finite spanning-tree gauge makes every link in a smaller
cube $C_Q\epsilon$-close to the identity. One direct justification is to
fill each finite tree cycle by finitely many plaquettes; its holonomy is a
product of conjugated plaquette holonomies. Unitary telescoping and
$\|I-U_p\|_{\rm HS}^2=2s_p$ bound its distance by $C_Q\epsilon$.
The cube size and all cycle lengths are fixed.

Call a conditioned bath good for a block if some vertex gauge makes all
bath links in the required finite neighborhood $\delta$-close to identity.
This is an event of the bath alone (minimize the maximal distance over the
compact finite vertex gauge group). The preceding full-configuration event
implies this bath event, even though the gauge witnessing that implication
can depend on the full configuration. Thus its marginal probability is at
least $1/2$ once $\beta\ge\beta_0(\delta)$.

Choose translated two-link blocks with centers in a fixed-spacing sublattice,
and keep only an interior packed set, so their entire determining
neighborhoods are disjoint. Increase the spacing to include every saved
pointing offset. There are at least $c_0n^4$ such blocks for $n\ge n_0$.
No plaquette contains two free links, even inside a block. After conditioning
on every other link, the Wilson law factors over all free links. On a good
block, after its bath gauge, either free link has conditional action
\[
 S_b(U)=6N-\Re\operatorname{tr}(UH_b),\qquad \|H_b-6I\|\le C\delta,
\]
up to an irrelevant bath constant. At $H_b=6I$ there is a unique minimum
at $I$, a strictly positive Hessian, and a strictly positive gap outside a
fixed small chart. Compactness, smooth dependence and a smaller $\delta$
preserve these facts uniformly, with the minimum at $x_b^*$ in the chart.
Subtracting the minimum and integrating local Gaussian upper/lower bounds
shows that the one-link normalizer is between $c\beta^{-d/2}$ and
$C\beta^{-d/2}$; the outside-chart mass is exponentially small. Haar
coordinate Jacobians are uniformly positive and bounded there.

In the actual coordinates
\[
 z=\sqrt\beta(u-u_*),\qquad w=\sqrt\beta(v-v_*),
\]
the conditional product density is therefore bounded below by a fixed
$m_0>0$ on a fixed product of unit balls, uniformly on all good baths and
$\beta\ge\beta_0$. This is a direct nonlinear density estimate. It does
not invoke asymptotic Gaussian convergence and does not drop linear terms
of the observable at the conditional minimum.

\section{Proof of the theorem and precise implications}
For either $G=G_H[1]$ or $G=D[1]$, locality and the separated block packing
give, conditionally on the bath $\mathfrak b$,
\[
 G=G_0(\mathfrak b)+\sum_jG_j(U_{j,-},U_{j,+};\mathfrak b).
\]
Each original force-square summand meets at most one block, so every
summand is included. Since no plaquette contains two free links, \eqref{lower}
has the conditional form
\[
 L=L_0(\mathfrak b)+\sum_j\{L_{j,-}(U_{j,-};\mathfrak b)
                            +L_{j,+}(U_{j,+};\mathfrak b)\}.
\]
No bound on these additive terms is needed. Conditional independence between
blocks implies additivity of their conditional variances.

For a good block set
$Q(z,w)=\beta G_j(u_*+\beta^{-1/2}z,v_*+\beta^{-1/2}w;\mathfrak b)$.
Its mixed derivative equals the corresponding unscaled mixed derivative
of $G_j$, so it is at least $\kappa$ by \eqref{jetD} or \eqref{jetH} and
continuity. Apply Lemma~\ref{mixed} with additive terms
$(\beta/\sigma)L_{j,-}$ and $(\beta/\sigma)L_{j,+}$. This yields
\[
 \Var(\sigma G_j-L_{j,-}-L_{j,+}\mid\mathfrak b)
 \ge c\sigma^2\beta^{-2}
\]
on each good bath. The law of total variance, the $1/2$ probability of a
good bath for each block (independence of those events is not needed), and
the packing count give
\[
 \Var(\sigma G-L)
 \ge c\sigma^2\beta^{-2}n^4
 =cc_\beta^{-2}a^{-4}n^4
 =c'T^4a^{-8}.
\]
This proves both uniform-test statements. For a physical real $f$ with a
plateau equal to one, put all blocks in an interior sub-box whose collar
contains the full determining neighborhoods and all original pointing
shifts. Every varying summand there has its actual test equal to one, so
the calculation is unchanged. The number of blocks is at least
$c|Q'|a^{-4}$. Terms outside these neighborhoods are bath terms and cannot
affect the conditional lower bound. This proves the plateau statements.

\paragraph{What is excluded.}
Identity and arbitrary weighted Wilson-action counterterms, even with
cutoff-divergent coefficients, cannot make the coincident variance of the
unrescaled H insertion or its specified ordinary remainder cutoff-uniform.
The conclusion is stronger than invariance of variance under centering.
If a multiplicative real normalization $\Gamma_a$ is also applied, the
bound is multiplied by $|\Gamma_a|^2$. A bounded coincident variance within
this restricted subtraction class would require $|\Gamma_a|=O(a^4)$ for a
fixed plateau. This is a necessary condition, not a proposed physical
normalization or a nondegeneracy result.

\paragraph{What is not excluded.}
The theorem does not cover all lower operators built from several
plaquettes, source-degree-two contact counterterms, general current
redefinitions, or renormalized separated correlations. Nor does divergence
of a bare coincident variance rule out a continuum distribution constructed
with an explicitly prescribed contact extension. A subtraction of a
mixed-source scalar changes that coincident source coefficient and must be
recorded; it is not an identity of the unchanged XP3 source ($A_{12}^{bare}=0$).
There is no assertion that the theorem disproves MT36, whose norm retains
lower/contact sectors.

\paragraph{Target for O2/O3.}
The actual ordinary remainder cannot be treated as already finite merely
because its exact Casimir companion was included. A useful next target is
an upper bound of order $a^{-4}N_f$ (with only explicitly justified coupling
factors), together with a separated-kernel estimate and a specified
source-degree-two local contact extension. The bound must transport through
the genuine mixed maps with the original tests, saved meshes, selectors,
current lifts, effective scores and common domains. The present theorem
proves none of those upper or transport estimates.

\paragraph{Absolute masses and domains.}
All ordered inequalities here concern real observables under the complete
positive Wilson probability on the full compact configuration space. There
are no auxiliary selectors or gauge-fixing determinants; local bath gauges
are merely Haar-preserving changes of variables. For the corresponding
raw connected numerator the positive variance is multiplied by the original
$Z_\beta$; an independent field-free factor $e^s$ multiplies its absolute
mass by $|e^s|$. No lower bound for $Z_\beta$ is inserted. No inequality is
asserted for an individual signed/complex selected history or a complex wall.

\paragraph{Manuscript interface.}
The inputs are AI1--3 (pp.~635--636), AI33--35d (pp.~642--644), MT1
(p.~681), XP2--4 (p.~719), and PD12.4--8,20--25,39--56
(pp.~1183,1186--1187,1190--1193). The estimates left open
are MT36 (p.~692) and PD12.58--63 (p.~1194). The exact terminal identity
$\Var(O_H)=\Var(r_H)+\E b_{HH}$ is preserved; lower bounds for its first
two variances alone imply no sign or matching estimate for $\E b_{HH}$.
\end{document}
